\section{Triangular transport and the physical filtration}
\label{sec:ope-filtration}
The distance estimate in \eqref{eq:ope-match-distance} concerns the
physical filtration. A bare coefficient of low engineering degree need
not stay in the corresponding coordinate plane under renormalization.
We first estimate its transported plane, then identify limits of that
plane by their correlations. The two parts use different information:
triangular estimates for the first, and the comparison of local
observables for the second.

\subsection{Expansion of the low coordinates}
Fix a degree $d$, its finite coefficient space $F$, and an integer
$0\leq e<d$. Each coordinate has an engineering weight in $\{0,\ldots,d\}$.
Write $\mu(v)=\sum_\alpha|v_\alpha|$ and let $P$ retain the coordinates
of weight at most $e$. For a block ratio $L\geq1$, let
\[
 (D_Lv)_\alpha=L^{d-\deg\alpha}v_\alpha.
\]
Thus $P$ commutes with $D_L$ and
\begin{equation}\label{eq:ope-dilation-split}
 \mu(PD_Lv)\geq L^{d-e}\mu(Pv),\qquad
 \mu((1-P)D_Lv)\leq L^{d-e-1}\mu(v).
\end{equation}
Both inequalities follow by summing the coordinate inequalities.

Here is the exact one-step input needed from the calibrated construction.
At consecutive levels $i+1,i$, the readouts of the same bare vector are
$a,b$. There is an invertible linear map $B$ preserving $\ker P$, with
\[
 \mu(Bv),\ \mu(B^{-1}v)\leq K\mu(v),\qquad
 \mu(b-BD_La)\leq\epsilon\mu(a).
\]
The last estimate follows in the construction from the stronger bound
with the nonconstant-coordinate part of $a$ on the right. Its numerical
conditions are
\begin{equation}\label{eq:ope-triangle-numbers}
 K\geq1,\quad s_0,s\geq64K^2,\quad s\leq L,\quad
 \epsilon\geq0,\quad 2K\epsilon(1+s)\leq1.
\end{equation}
Here $s_0=\zeta_i$ and $s=\zeta_{i+1}$ are the column allowances.
These triangular estimates concern the actual readouts, including their
nonlinear corrections. Their production is still a construction input.

Since $B$ maps the finite-dimensional space $\ker P$ injectively into
itself, it maps it onto itself. Its inverse therefore also preserves
$\ker P$. Applying this observation to $Bv-PBv$ gives
\begin{equation}\label{eq:ope-low-inverse}
 \mu(Pv)=\mu(PB^{-1}PBv)\leq K\mu(PBv).
\end{equation}
Suppose now that $\mu(a)\leq s\mu(Pa)$, and set $Y=\mu(PD_La)$.
By \eqref{eq:ope-dilation-split} and $s\leq L$,
\[
 \mu(D_La)\leq Y+L^{d-e-1}s\mu(Pa)\leq2Y,
 \qquad \mu(a)\leq sY.
\]
The error $E=\mu(b-BD_La)$ consequently satisfies
$2KE\leq Y$. On the other hand, \eqref{eq:ope-low-inverse} gives
$Y\leq K(\mu(Pb)+E)$, so $Y\leq2K\mu(Pb)$. It follows that
\begin{equation}\label{eq:ope-cone-step}
 \mu(b)\leq3KY\leq6K^2\mu(Pb)\leq s_0\mu(Pb),\qquad
 L^{d-e}\mu(Pa)\leq s\mu(Pb).
\end{equation}
If instead $Pa=0$, \eqref{eq:ope-dilation-split}, $\epsilon\leq1$,
and $\mu(a)\leq\mu(D_La)$ give
\begin{equation}\label{eq:ope-transverse-step}
 \mu(b)\leq(K+1)\mu(D_La)
       \leq sL^{d-e-1}\mu(a).
\end{equation}
The lost power of $L$ distinguishes the transverse estimate from the
expansion of the low coordinates.

\subsection{Subtracting the transported low plane}
Fix a finest level $N$, a period, and a reference level $U\leq N$.
Let $T_k:F\to F$ be the linear readout at level $k$, for $U\leq k\leq N$.
Thus $T_N=1$ and $T_U=K_N$, where $K_N$ is the invertible map to
reference coefficients. For each adjacent pair assume the preceding
one-step estimates, with $a_k=L_ka_{k+1}$ for the meshes. Define
\[
 J_k=\prod_{i<k}\zeta_{i+1},\qquad
 H_k=\prod_{i<k}(1+\zeta_i)\zeta_{i+1}.
\]
All these factors are at least one. If $Pv=v$, backward induction in
\eqref{eq:ope-cone-step} gives
\begin{equation}\label{eq:ope-plane-expansion}
 \mu(T_kv)\leq\zeta_k\mu(PT_kv),\qquad
 J_k(a_k/a_N)^{d-e}\mu(v)\leq J_N\mu(PT_kv).
\end{equation}
For the second inequality, multiply the one-step lower expansion by
$J_i$, use $J_{i+1}=\zeta_{i+1}J_i$, and telescope the mesh ratios.
The terminal equality $T_Nv=v$ starts both inductions.

In particular $PT_k$ restricted to $\operatorname{ran}P$ is injective.
Its domain and codomain have equal finite dimension, so it is onto.
For any $y\in F$ choose $v\in\operatorname{ran}P$ with
$PT_kv=PT_ky$. The first bound in \eqref{eq:ope-plane-expansion} shows
\begin{equation}\label{eq:ope-plane-cancel}
 PT_k(y-v)=0,\qquad
 \mu(T_k(y-v))\leq(1+\zeta_k)\mu(T_ky).
\end{equation}
This cancellation must be repeated after each step: the fixed coordinate
complement is not asserted to be invariant under the exact readouts.
If $PT_{i+1}y=0$, first apply \eqref{eq:ope-transverse-step}, then
\eqref{eq:ope-plane-cancel} at level $i$. The cost is
$(1+\zeta_i)\zeta_{i+1}L_i^{d-e-1}$.
Induction therefore gives, whenever $U\leq k\leq j\leq N$ and
$PT_jy=0$, a $v\in\operatorname{ran}P$ such that
\[
 PT_k(y-v)=0,\qquad
 H_k\mu(T_k(y-v))\leq
 H_j(a_k/a_j)^{d-e-1}\mu(T_jy).
\]
At an induction step add the newly subtracted low vector to the old
one; their sum remains in $\operatorname{ran}P$. The recurrence for
$H_k$ and the identity $a_i=L_ia_{i+1}$ give the displayed product.

Cancel first at level $k$, then transport to $U$, and discard the
factor $H_U\geq1$. For every reference vector $x$ this yields a vector
$y$ in the transported plane $K_N\operatorname{ran}P$ with
\begin{equation}\label{eq:ope-transport-distance}
 \|x-y\|\leq F_{k,d,e}\,q_k(x),\qquad
 F_{k,d,e}=H_k(1+\zeta_k)(a_U/a_k)^{d-e-1},
\end{equation}
where $q_k(x)=\mu(T_kK_N^{-1}x)$, and the reference norm is chosen so
that $\|z\|\leq\mu(z)$. Indeed apply the preceding result to
$K_N^{-1}x$ and use $T_U=K_N$.
For comparison, the calibrated bound
$q_i(x)\leq\zeta_{i+1}L_i^d q_{i+1}(x)$ gives by forward induction
\begin{equation}\label{eq:ope-whole-distance}
 \|x\|\leq J_k(a_U/a_k)^d q_k(x).
\end{equation}
At $U$ use $q_U(x)=\mu(x)$ and $J_U\geq1$; the recurrence for $J$
and the mesh identity give each subsequent step.

\subsection{Subpower calibration and smeared discrepancies}
\label{ex-column-subpower}
The cumulative column factor used above is subpower for a precise
reason. Put $s_i=\log(1/a_i)\geq0$ and
$C_N=\sum_{i<N}\log\zeta_{i+1}=\log J_N$.
The schedule provides numbers $h_i\to\infty$ and the bound
\[
 0\leq\log\zeta_{i+1}\leq h_i^{-2}\log L_i,
 \qquad s_{i+1}-s_i=\log L_i>0.
\]
The first inequality follows also when the scheduled logarithmic cost
contains other nonnegative charges, since it then dominates
$\log\zeta_{i+1}$. Given $\epsilon>0$, choose $k$ so that
$h_i^{-2}\leq\epsilon/2$ for $i\geq k$. Telescoping gives
\[
 C_N\leq C_k+\tfrac\epsilon2(s_N-s_k)\leq C_k+\tfrac\epsilon2s_N.
\]
Eventually $C_k\leq\epsilon s_N/2$, because $s_N\to\infty$.
Thus $C_N=o(s_N)$. In particular $J_N$ has logarithmic rate zero;
so does $\zeta_N$, since $1\leq\zeta_N\leq J_N$ for $N\geq1$.

The calibrated reverse estimate is
\[
 \mu(T_{i+1}K_N^{-1}x)\leq\zeta_{i+1}\mu(T_iK_N^{-1}x).
\]
Starting at $T_UK_N^{-1}x=x$ and enlarging the product to the full
prefix gives
$\mu(T_iK_N^{-1}x)\leq J_i\mu(x)$.
At $i=N$, therefore,
\[
 \|K_N^{-1}x\|\leq J_N\mu(x)
            \leq J_N\,\dim(F)\,\|x\|,
\]
using the coordinate supremum norm for the last inequality. For every
fixed radius $R$ and $\epsilon>0$, both $J_N$ and
$\dim(F)R$ are eventually bounded by $e^{\epsilon s_N/2}$.
Consequently
\begin{equation}\label{eq:ope-bounded-bare}
 \sup_p\sup_{\|x\|\leq R}\|K_{N,p}^{-1}x\|
       \leq e^{\epsilon s_N}
 \quad\hbox{eventually in }N.
\end{equation}
The same estimate applies to any bounded family along the ordered
cutoff filter. The calibration inequalities are uniform in the period;
no new period extraction is needed.

We will also use the following passage from separated lattice bounds
to smeared tests. Let a trial $Z_{N,p}$ replace one distinguished field
in an $n$-point correlation, and suppose that, for each fixed spectator
list and every $\delta>0$, its correlation has absolute value at most
$A_\delta|r_{N,p}|$ whenever all physical pairwise separations are at
least $\delta$. Assume $r_{N,p}\to0$ along the ordered cutoff filter.
For every smooth test $\varphi$ compactly supported in the separated
configuration space, the corresponding lattice sum, with the factor
$a_N^{4n}$, tends to zero.

Indeed the test support $K$ is compact and has positive distance from
every pairwise diagonal. Choose a smaller clearance $\delta>0$.
For sufficiently large periods, its lattice representatives have that
clearance also in the periodic metric: a bounded set does not meet a
periodic translate of itself near a diagonal. Express the weighted
lattice sum exactly as the integral of its piecewise constant function
on the mesh cells in the fundamental region. A nonzero cell value has
its corner in $K$. The cell diameter tends to zero, so every such cell
lies in a fixed compact enlargement $B$ of $K$. The absolute value of
the step function is at most
\[
 \boldsymbol1_B\,\|\varphi\|_\infty A_\delta |r_{N,p}|.
\]
Its integral is bounded by
$|B|\|\varphi\|_\infty A_\delta |r_{N,p}|$, which tends to zero.
This is a norm estimate on the same ordered filter, not an interchange
of mesh and volume limits. Coordinate ties do not cause a problem:
the test support is separated in the Euclidean sense, and one
separating coordinate for each pair suffices to control its periodic
separation.

\subsection{Identifying the limiting plane by tests}
We spell out the finite-dimensional limit argument, since convergence
of the moving planes themselves is neither needed nor asserted.
Let $V$ be a finite-dimensional complex normed space, $p_c$ projections
converging in operator norm to a projection $p$, and $f_{\theta,c}$
linear functionals converging in norm to $f_\theta$, one for each test
$\theta$. Assume
\begin{equation}\label{eq:ope-test-kernel}
 \ker p=\bigcap_\theta\ker f_\theta.
\end{equation}
The index $c$ tends along any nontrivial filter. Suppose vectors $u_c$
satisfy, for every $\theta$,
\begin{equation}\label{eq:ope-normalized-tests}
 f_{\theta,c}(u_c)\longrightarrow M_\theta,\qquad
 \frac{f_{\theta,c}(u_c-p_cu_c)}{1+\|p_cu_c\|}\longrightarrow0.
\end{equation}
Then there is $u\in\operatorname{ran}p$ with
$f_\theta(u)=M_\theta$ for all $\theta$.

To prove this, put $b_c=(1+\|p_cu_c\|)^{-1}$ and $z_c=b_cp_cu_c$.
The pairs $(b_c,z_c)$ lie in the compact set $[0,1]\times\overline B_1$.
Along a finer nontrivial filter they converge to $(b,z)$, with
$b+\|z\|=1$. Since $p_cz_c=z_c$, convergence of the projectors gives
$pz=z$. Equations \eqref{eq:ope-normalized-tests} imply
$f_\theta(z)=bM_\theta$. If $b=0$, \eqref{eq:ope-test-kernel} forces
$pz=0$, hence $z=0$, contradicting $b+\|z\|=1$. Thus $b>0$ and
$u=b^{-1}z$ proves the claim. There is only one compact extraction;
the limits for every test pass to its refinement. No extraction indexed
by the tests is used.

Here is how the required hypotheses arise for the transported plane.
We make explicit the remaining comparison input. A local trial at degree
$e$ is a family of observables $Z_N$, indexed along the ordered cutoff
filter, with the following properties on the late tail of that filter.
(i) Locality: there is one fixed finite stencil $X$ such that each $Z_N$
is a measurable gauge-invariant function of the link variables in the
translate of $X$ at its lattice origin, and the observable at any other
origin is the literal translate of the one at the origin. (ii) Eventual
continuity: $Z_N$ is continuous for all sufficiently late cutoffs.
(iii) Two prices: there are nonnegative sequences $c_N$ and $b_N$ of zero
logarithmic rate such that $a_N^e|Z_N|\le b_N$ everywhere, and such that
$a_N^eZ_N$ satisfies the prescribed left-derivative bounds with constant
$c_N$ on the healthy configurations, those at which no site of the
stencil is incident to a plaquette whose chord distance from the identity
exceeds the health threshold of the cutoff; no derivative
bound is required on unhealthy configurations. These bounds are uniform in
the lattice origin. Here the derivative
order, spatial weight power and direction threshold are the purchased
cutoff-dependent parameters: the bound holds simultaneously for every
ordered left derivative of length at most that cutoff's purchased order,
with directions satisfying its weighted threshold. Bounds only for each
fixed derivative order would not suffice. The local comparison
construction supplies degree-$e$ labels $u_c$ with subpower bare norm,
whose difference from that trial has a centered price tending to zero.
It also supplies the projections $p_c\to p$ of
\eqref{eq:ope-test-kernel}. Labels in $\ker p_c$ with subpower bare norm
have vanishing tested correlations. Bounded reference labels have
subpower bare norm by \eqref{eq:ope-bounded-bare}. The local comparison
and projection assertions are substantive inputs still to be proved
from the construction. All statements use
the same ordered cutoff filter.

Section~\ref{sec:finite-rank-kernel-transport} supplies projections of
this kind: its idempotents $A_c$ converge to $T$ with
$\ker T=\bigcap_\theta\ker f_\theta$, by
\eqref{eq:finite-rank-exact-kernel}, so one may take $p_c=A_c$ and
$p=T$ in \eqref{eq:ope-test-kernel}. Its vanishing statement
\eqref{eq:finite-rank-silent} is displayed for bounded families. The
argument there applies verbatim to families of subpower bare
(inverse-calibrated) norm, since its local equation input is stated for
such families and boundedness is used only to enter that input. The use
here requires that subpower form, for the same projections, tests and
filter. That section in turn rests on its own inputs: its local
equation input, the purchased comparison
\eqref{eq:finite-rank-purchased-matrix} and the equation count
\eqref{eq:finite-rank-equation-count}. The local trial comparison above
is a separate input.

For clarity, these inputs imply the second limit in
\eqref{eq:ope-normalized-tests} without assuming $u_c$ bounded.
Multiplication by $b_c=(1+\|p_cu_c\|)^{-1}\leq1$ preserves subpower
bare norm. The vectors $b_cp_cu_c$ are bounded, so also have subpower
bare norm. Their difference $b_c(u_c-p_cu_c)$ consequently has
subpower bare norm and belongs to $\ker p_c$, because $p_c^2=p_c$.
The vanishing-test assertion gives exactly the normalized residual
limit. The first limit follows by comparing the trial with $u_c$:
the centered estimate gives a separated correlation bound by a
constant times the price, and the step-integral estimate just proved
sends this bound to zero in every smeared correlation. Linearity here is legitimate at finite cutoff,
since the trials are bounded measurable functions and the remaining
fields are continuous on the compact configuration space.

It remains to check the degree of a bounded transported low vector
$v_c\in K_N\operatorname{ran}P$. Write $x_c=K_N^{-1}v_c$.
Taking $k=U$ in \eqref{eq:ope-plane-expansion} gives
\[
 a_N^{-(d-e)}\|x_c\|
 \leq a_U^{-(d-e)}J_N\mu(v_c).
\]
Section~\ref{ex-column-subpower} gives zero logarithmic rate for $J_N$. Thus the bare
coefficients of $w_c=a_N^{-(d-e)}v_c$ have subpower norm.
The calibrated local-section bounds give subpower prices for
$a_N^d\mathcal O_d(w_c)$. By linearity,
\[
 a_N^d\mathcal O_d(w_c)=a_N^e\mathcal O_d(v_c).
\]
The same literal observable, with its original finite support, is
therefore a degree-$e$ trial in the sense above, provided the calibrated
local-section bounds give for it the two prices of (iii) on the healthy
configurations, together with literal translation and eventual
continuity (i)--(ii). We use those bounds in exactly this form; their
proof is part of the open local comparison input stated above. If $v_c\to v$ along a
refinement, its tested correlations converge to those of the reference
label $v$: convergence of linear tests on a fixed basis is convergence
in operator norm. Apply the preceding comparison and compact argument
to obtain a degree-$e$ label having all the same smeared correlations.
The smooth correlation kernels consequently agree: their difference
integrates to zero against every compactly supported smooth test on
the separated configuration domain, hence vanishes there. By the
physical null-space definition, the class of $v$ belongs to $E_e$.

Only boundedness and low-plane membership on the refined filter are
needed. To apply comparison inputs stated on the original filter,
replace $v_c$ by zero whenever either $\|v_c\|>\|v\|+1$ or low-plane
membership fails. This modified family is bounded and in the plane
on the original filter, and equals $v_c$ eventually on the refinement.
It has the same limit and the same limiting correlations there.

\subsection{The distance estimate after extraction}
Suppose reference labels $v_c\to v$ on a refinement of the ordered
cutoff filter and $q_k(v_c)\leq V$ eventually, with $k$ fixed.
Choose $y_c\in K_N\operatorname{ran}P$ by
\eqref{eq:ope-transport-distance}. They are eventually bounded, because
$\|y_c\|\leq\|v_c\|+F_{k,d,e}V$. Extract a further nontrivial filter
on which $y_c\to y$. The preceding argument puts $[y]$ in $E_e$,
and continuity of the norm gives
\begin{equation}\label{eq:ope-physical-distance}
 [y]\in E_e,\qquad \|v-y\|\leq F_{k,d,e}V.
\end{equation}
The whole-vector bound \eqref{eq:ope-whole-distance} passes to the
limit directly. Both estimates hold for the same limit $v$; only the
auxiliary approximant $y$ requires a further extraction.

Finally these finite-scale factors have the advertised rates. By Section~\ref{ex-column-subpower},
$\log J_k=o(s_k)$ and $\log\zeta_k=o(s_k)$, where
$s_k=\log(1/a_k)$. The column allowances also satisfy $\zeta_k\geq2$.
In fact
\[
 H_k\zeta_k\leq J_k^3\zeta_0.
\]
This holds at $k=0$. For the induction step use
$1+\zeta_k\leq\zeta_k\zeta_{k+1}$ and the recurrences for $H,J$:
\[
 H_{k+1}\zeta_{k+1}
 \leq H_k\zeta_k\zeta_{k+1}^3
 \leq J_k^3\zeta_0\zeta_{k+1}^3
 =J_{k+1}^3\zeta_0.
\]
Thus $H_k$ also has zero logarithmic rate. If $k_n\to\infty$, $t_n\to+\infty$, and $s_{k_n}/t_n\to a$,
then $F_{k_n,d,e}$ has rate $(d-e-1)a$, while the factor in
\eqref{eq:ope-whole-distance} has rate $da$.
Taking $d=q$ and $e=m-1$ proves precisely the factors in
\eqref{eq:ope-match-norm}--\eqref{eq:ope-match-distance}, once a
cutoff match with quota $V_n$ and convergent reference labels has been
produced. This argument preserves the exact reference labels and
physical filtration; it does not require convergence of the individual
renormalization matrices.
